\section*{Hoofdstuk \ref{ch:b2:climate-variability} --- Klimaatvariabiliteit}

\begin{solution}{exo:b2:climate-variability:1}
(a) $29.27 - 26.71 = \qty{2.56}{\celsius}$ in 2026; $28.79 - 26.73 = \qty{2.06}{\celsius}$ in 1997.
(b) September 2026 was zowel de warmere maand, met \qty{0.48}{\celsius}, als de grotere anomalie, met
\qty{0.50}{\celsius}. (c) Elk jaar wordt vergeleken met een basisperiode van 30 jaar eromheen, die
om de vijf jaar wordt verschoven; de septemberklimatologie van de jaren 1980 en 1990 verschilt een beetje van die van
de jaren 2010 en 2020.
\end{solution}

\begin{solution}{exo:b2:climate-variability:2}
(a) JJA: $(0.79 + 0.95 + 1.28)/3 = 1.007$; JAS: $(0.95 + 1.28 + 1.55)/3 = 1.260$; ASO: $(1.28 +
1.55 + 1.68)/3 = 1.503$; SON: $(1.55 + 1.68 + 2.01)/3 = 1.747$; OND: $(1.68 + 2.01 + 2.02)/3 =
1.903$, alle in \unit{\celsius}. (b) Afgerond op één decimaal: 1.0, 1.3, 1.5, 1.7, 1.9, de waarden van
de tabel. (c) In ASO 2023, met 1.5.
\end{solution}

\begin{solution}{exo:b2:climate-variability:3}
(a) De index volgt de afwijking van de druk in Tahiti min die in Darwin: onder het gemiddelde
in Tahiti en erboven in Darwin is hij negatief. (b) El~Niño. (c) Zwakkere oostelijke passaatwinden, zelfs
westenwinden dicht bij het oppervlak in het westen, en minder regen dan gewoonlijk boven Indonesië, omdat de
diepe convectie naar het oosten is verschoven.
\end{solution}

\begin{solution}{exo:b2:climate-variability:4}
(a) $804.1/880.6 = \qty{91.3}{\percent}$ van het gemiddelde: binnen 90 tot \qty{110}{\percent}, een
normaal seizoen, al in het onderste deel. (b) Het seizoen komt in totaal op $556 + 659 + 427 + 252 =
\qty{1894}{mm}$; juli brengt daarvan $659/1894 = \qty{34.8}{\percent}$, meer dan een derde.
\end{solution}

\begin{solution}{exo:b2:climate-variability:5}
(a) $\tau = 1025 \times 3992 \times 50/16.1 = \qty{1.271e7}{s} = \qty{147.1}{\day}$ voor
\qty{50}{m}, en het dubbele, \qty{294.2}{\day}, voor \qty{100}{m}. (b) $\mathrm{e}^{-30.44/147.1} =
0.813$ en $\mathrm{e}^{-30.44/294.2} = 0.902$. (c) De dempingstijd is evenredig met $h$: een
ondiepere menglaag slaat bij dezelfde flux minder warmte op en vergeet een anomalie sneller, dus het
geheugen van de oceaan, en de voorspelbaarheid die het geeft, nemen af.
\end{solution}

\begin{solution}{exo:b2:climate-variability:6}
(a) Drie maanden zijn 91.3 dagen: $1.2 \times \mathrm{e}^{-91.3/147.1} = \qty{0.65}{K}$; zes maanden
zijn 182.6 dagen: $1.2 \times \mathrm{e}^{-182.6/147.1} = \qty{0.35}{K}$. (b) De anomalie dooft
traag uit, dus ze is een betere gok voor de anomalie van de volgende maand dan nul. Maar de willekeurige forcering voegt
elke dag nieuwe anomalieën toe, onbekend op het moment van de verwachting; terwijl de beginanomalie uitdooft,
nemen die een steeds groter aandeel in, en drijft de verwachting af naar de \omterm{def:b2:climate-variability:modes}{klimatologie}, waar ze niet
langer helpt.
\end{solution}

\begin{solution}{exo:b2:climate-variability:7}
(a) $160 \times 86\,400 = \qty{1.38e7}{J}$, \qty{13.8}{MJ} per vierkante meter en per dag. (b) De
laag bevat $(1000 - 850) \times 100/9.81 = \qty{1529}{kg}$ lucht per vierkante meter, met een
warmtecapaciteit van $1529 \times 1004.5 = \qty{1.536e6}{J.m^{-2}.K^{-1}}$; ze zou $1.38 \times
10^7/1.536 \times 10^6 = \qty{9.0}{K}$ per dag opwarmen. (c) De laag straalt warmte uit in het infrarood, naar
de ruimte en naar de lucht erboven; en de warme lucht stijgt op en mengt zich met de laag erboven, of wordt
door de winden afgevoerd, waaronder de instroming van koelere lucht die de hittedepressie zelf aanzuigt.
\end{solution}

\begin{solution}{exo:b2:climate-variability:8}
(a) $(-1.93 - 1.11 - 1.98)/3 = -1.673$, ongeveer $-1.67$. (b) Het is de laagste winter van de reeks.
(c) Een sterk negatieve NAO: een zwakke IJslandse depressie en een zwak Azorenhoog, een zwakkere straalstroom en een stormbaan
zuidelijker dan gewoonlijk; een koude, droge winter in Noord-Europa, een natte in Zuid-Europa,
en meer uitbraken van koude lucht in het oosten van de Verenigde Staten.
\end{solution}

\begin{solution}{exo:b2:climate-variability:9}
(a) $9.3 - 0.4 = \qty{8.9}{mm}$ per dag, meer dan 2.5: in het domein. (b) Op het zuidelijk halfrond
is de zomer december--februari: $6.1 - 4.2 = \qty{1.9}{mm}$ per dag, minder dan 2.5: erbuiten,
hoewel het er meer regent dan op de meeste plaatsen van het domein. (c) $1894/122 = \qty{15.5}{mm}$ per
dag; Kerala zou er alleen buiten liggen als zijn winterhoeveelheid $15.5 - 2.5 = \qty{13.0}{mm}$ per dag overschreed.
(d) Een \omterm{def:b2:climate-variability:monsoon}{moesson} is een omkering met de seizoenen: een natte plaats die het hele jaar door nat is, zoals het tweede
station, heeft geen \omterm{def:b2:climate-variability:monsoon}{moesson}, en een droge plaats met een kort, hevig regenseizoen heeft er wel een.
\end{solution}

\begin{solution}{exo:b2:climate-variability:10}
(a) Differentiëren geeft $\ddot T = a\dot H = -abT$: de vergelijking van een harmonische oscillator met
hoekfrequentie $\omega = \sqrt{ab}$ en periode $2\pi/\sqrt{ab}$. (b) $2\pi/\sqrt{ab} = 4$ geeft $ab =
(2\pi/4)^2 = \qty{2.47}{yr^{-2}}$. (c) Met $T = u\,\mathrm{e}^{-rt}$, $H = v\,\mathrm{e}^{-rt}$ is
$\dot T = (\dot u - ru)\,\mathrm{e}^{-rt} = (av - ru)\,\mathrm{e}^{-rt}$, dus $\dot u = av$, en
evenzo $\dot v = -bu$: $(u, v)$ oscilleert zonder demping, en $(T, H)$ is die oscillatie
vermenigvuldigd met $\mathrm{e}^{-rt}$. In één periode daalt de amplitude met een factor $\mathrm{e}^{-0.25 \times
4} = \mathrm{e}^{-1} = 0.37$. (d) Een forcering: de willekeurige schommelingen van de passaatwinden, die
het systeem telkens opnieuw een duw geven, en de \omterm{def:b2:climate-variability:bjerknes}{Bjerknes-terugkoppeling}, die de demping zwakker of
zelfs negatief maakt.
\end{solution}

\begin{solution}{exo:b2:climate-variability:11}
(a) $A = R^2 \,\Delta\lambda\,(\sin\ang{5} - \sin(\ang{-5})) = (6.371 \times 10^6)^2 \times 0.8727
\times 0.1743 = \qty{6.17e12}{m^2}$. (b) $Q = 1025 \times 3992 \times 100 \times 2 \times 6.17
\times 10^{12} = \qty{5.05e21}{J}$. (c) Zes maanden zijn $182.6 \times 86\,400 = \qty{1.578e7}{s}$;
per vierkante meter $8.18 \times 10^8/1.578 \times 10^7 = \qty{51.9}{W/m^2}$, een half jaar lang
volgehouden. (d) Uit de oceaan zelf: zwakkere passaten brengen minder koud water omhoog in het oosten, en het warme
water van het westen, opgeslagen boven de diepere thermocline, verspreidt zich naar het oosten; een groot deel van de opwarming
kan komen van warmte die al in de oceaan zat, verplaatst in plaats van gewonnen.
\end{solution}

\begin{solution}{exo:b2:climate-variability:12}
(a) De lus geeft $0.3g$ terug, dan $0.3g^2$, enzovoort: het totaal is $0.3(1 + g + g^2 + \dots) =
0.3/(1 - g)$, een meetkundige reeks die convergeert voor $g < 1$: \qty{0.6}{K}, \qty{1.5}{K} en
\qty{3.0}{K} voor $g = 0.5$, 0.8 en 0.9. (b) De reeks divergeert: de anomalie zou onbeperkt
groeien; de groei moet worden gestopt door processen die de lus weglaat, de demping door de warmteflux
aan het oppervlak en het leeglopen van de warmwatervoorraad dat de opladingsoscillator beschrijft. (c)
Hoe warmer het zeeoppervlak, hoe gemakkelijker een kleine opwarming diepe convectie en regen uitlokt,
waarvan de verwarming de winden aandrijft; dat is het seizoen waarin de effecten van ENSO het sterkst zijn.
\end{solution}

\begin{solution}{pb:b2:climate-variability:1}
\textbf{1.} $(0.26 + 0.75 + 1.13)/3 = \qty{0.713}{\celsius}$.
\textbf{2.} MJJ $(0.75 + 1.13 + 1.52)/3 = 1.133$; JJA $(1.13 + 1.52 + 1.80)/3 = 1.483$; JAS $(1.52 +
1.80 + 2.06)/3 = 1.793$.
\textbf{3.} ASO $(1.80 + 2.06 + 2.26)/3 = 2.040$; SON $(2.06 + 2.26 + 2.42)/3 = 2.247$; OND $(2.26
+ 2.42 + 2.36)/3 = 2.347$; NDJ $(2.42 + 2.36 + 2.33)/3 = 2.370$.
\textbf{4.} Afgerond op één decimaal 0.7, 1.1, 1.5, 1.8, 2.0, 2.2, 2.3 en 2.4: de tabel.
\textbf{5.} Eén enkele maand bevat weersruis; ENSO duurt vele maanden, en het gemiddelde van drie
verwijdert veel van de ruis terwijl het trage signaal behouden blijft.
\textbf{6.} In JJA 1997, met 1.483, afgerond 1.5; de piek was 2.4 in NDJ.
\textbf{7.} Van ASO 2018 tot MJJ 2019, en van SON 2019 tot FMA 2020.
\textbf{8.} Tien seizoenen, en zes.
\textbf{9.} In ASO 2020; binnen 2020 liggen maar vier seizoenen, ASO tot NDJ, op of onder $-0.5$, één
te weinig voor de regel, en DJF 2021 maakt hem compleet.
\textbf{10.} Eén El~Niño, SON 2018 tot MAM 2019, zeven seizoenen (AMJ 2019 zakt naar 0.37); de
La~Niña begint in AMJ 2020 ($-0.55$), vier seizoenen vóór het moment waarop de ONI haar plaatst, in ASO.
\textbf{11.} Nee: de relatieve index komt eind 2019 nooit boven 0.29. De hele tropische oceaan was
toen ongeveer een halve graad warmer dan zijn basisperiode; het Niño-3.4-gebied was warm in
absolute zin, maar nauwelijks warmer dan de tropen eromheen.
\textbf{12.} Elf.
\textbf{13.} 8, 7, 10, 4, 5, 6, 12, 5, 9 en 3 jaar.
\textbf{14.} De kortste is 3 jaar (2023 tot 2026), de langste 12 (1997 tot 2009).
\textbf{15.} $69/10 = 6.9$ jaar op beide manieren: de som van opeenvolgende tussenpozen is de tijd van het
eerste begin tot het laatste, $2026 - 1957 = 69$ jaar.
\textbf{16.} Ze heeft al vijf seizoenen geduurd, MAM tot JAS, op of boven \qty{0.5}{\celsius}, en
piekte op 2.2, boven 1.5; haar begin, MAM 2026, is het eerste seizoen van de reeks, en latere seizoenen
kunnen de reeks alleen verlengen, niet haar begin verschuiven.
\textbf{17.} Nee, en dat hoeft ook niet: de 25 episodes van alle sterkten komen ongeveer om de drie jaar, zoals de FAQ zegt;
de sterke zijn een deelverzameling, dus ze zijn zeldzamer, en hun tussenpozen lopen van 3 tot 12 jaar.
\textbf{18.} $\tau = 1025 \times 3992 \times 50/16.1 = \qty{1.27e7}{s} = \qty{147}{\day}$.
\textbf{19.} Menglaag: $1025 \times 3992 \times 50 = \qty{2.05e8}{J.m^{-2}.K^{-1}}$; luchtkolom:
$(10^5/9.81) \times 1004.5 = 10\,194 \times 1004.5 = \qty{1.02e7}{J.m^{-2}.K^{-1}}$, 20 keer
minder.
\textbf{20.} $1.024 \times 10^7/(1025 \times 3992) = \qty{2.5}{m}$ zeewater.
\textbf{21.} $1.024 \times 10^7/16.1 = \qty{6.36e5}{s} = 7.4$~dagen.
\textbf{22.} Een anomalie van de lucht alleen wordt in ongeveer een week gedempt, een van de menglaag in
ongeveer vijf maanden, en de warme laag boven de thermocline is nog dieper: alleen de oceaan houdt
lang genoeg warmte vast om een anomalie van het ene seizoen naar het volgende te dragen, terwijl de atmosfeer
de winden levert die de twee koppelen.
\textbf{23.} Sinds 1950 zijn sterke El~Niño-episodes gemiddeld om de \textbf{6.9 jaar} begonnen,
ongeveer zeven.
\end{solution}
